% 3.3 对称点和对称线的分类
\section{对称点和对称线的分类}

对某些格子，布里渊区的形状是唯一的；而对另一些格子，则有两种或更多种可能的形状，取决于基本平移的相对大小和它们之间的夹角。正空间 14 种布拉维格子的各种布里渊区示于图 3.2--3.15。

若 $(\mathbf{k}_{1} - \mathbf{k}_{2}) = a_{1}\mathbf{g}_{1} + a_{2}\mathbf{g}_{2} + a_{3}\mathbf{g}_{3}$，其中 $a_{1}$、$a_{2}$、$a_{3}$ 为整数，则称矢量 $\mathbf{k}_{1}$ 与 $\mathbf{k}_{2}$ \textit{等价}。由于其构造方式，布里渊区内部任何两点都不等价；但区面上的一个点可以与区面上一个或多个其他点等价。例如，一个面的中点总与对面中点等价，因为相对的面总相差一个倒格子矢量。

给定一个布里渊区和区内一点 $\mathbf{k}$，相应晶系的全对称点群 $\mathbf{P}$ 中存在某些元素，它们把 $\mathbf{k}$ 变换到它自身或某个与它等价的 $\mathbf{k}$ 矢量。这些元素构成 $\mathbf{P}$ 的一个子群，记作 $\mathbf{P}(\mathbf{k})$，称为 $\mathbf{k}$ 的\textit{对称群}（symmetry group）。现在可以给出对称点、对称线和对称面的定义。

\noindent\textbf{定义 3.3.1}\quad 若存在 $\mathbf{k}$ 的一个邻域 $N$，使得在 $N$ 内除 $\mathbf{k}$ 以外没有任何点具有对称群 $\mathbf{P}(\mathbf{k})$，则称 $\mathbf{k}$ 是一个\textit{对称点}（point of symmetry）。

也就是说，若 $\mathbf{k}' \in N$ 且 $\mathbf{k}' \neq \mathbf{k}$，则 $\mathbf{P}(\mathbf{k}')$ 总是 $\mathbf{P}(\mathbf{k})$ 的真子群。我们可以粗略地说，点 $\mathbf{k}$ 比它周围所有的点对称性都高。

\noindent\textbf{定义 3.3.2}\quad 若在 $\mathbf{k}$ 的任一充分小邻域 $N$ 中，总有一条线（一个平面）穿过 $\mathbf{k}$，其上所有点都与 $\mathbf{k}$ 具有相同的对称群，则称 $\mathbf{k}$ 是一条\textit{对称线}（line of symmetry）（一个\textit{对称面}（plane of symmetry））。

最后，若在 $\mathbf{k}$ 的任一充分小邻域 $N$ 中所有点都具有相同的对称群，则称 $\mathbf{k}$ 为\textit{一般点}（general point）。对一般点，$\mathbf{P}(\mathbf{k})$ 仅由恒等元构成。对对称面，$\mathbf{P}(\mathbf{k})$ 由恒等元和一个镜面反射构成，只有两个元素。

\noindent\textbf{定义 3.3.3}\quad 对每个布里渊区存在一个\textit{基本域}（basic domain）$\Omega$，使得 $(\Sigma_{R} R\Omega)$ 等于整个布里渊区，其中 $R$ 取遍相关晶系全对称点群 $\mathbf{P}$ 的元素。

因此
\begin{equation*}
\Omega = \frac{8\pi^{3}}{V|\mathbf{P}|}
\tag{3.3.1}
\end{equation*}
其中 $|\mathbf{P}|$ 是 $\mathbf{P}$ 的阶。

对称点的求法是：对交于该点的各个平面使用式 (3.2.6)，算出每个面心和每个顶点的坐标，然后看在相关全对称点群 $\mathbf{P}$ 的操作下（见表 3.1），这些点是否具有邻域内其他点都不具备的额外对称操作。这里所说的对称操作是指把矢量 $\mathbf{k}$ 变换为 $\mathbf{k}$ 本身或与 $\mathbf{k}$ 相差一个倒格子矢量的矢量 $\mathbf{k}'$ 的操作。类似的手续也可用于辨认对称线。所有布拉维格子布里渊区的各种对称点和对称线都已在图 3.2--3.15 的图中标出。我们在图 3.2--3.15 中用于对称点和对称线的字母几乎都是标准的，并尽可能与前人（Bouckaert, Smoluchowski, and Wigner 1936, Herring 1942, Koster 1955, Luehrmann 1968, Meijer and Bauer 1962, Miller and Love 1967）的记号一致。

% ---- 图 3.2 -- 3.6（MinerU 提取图；题注坐标已逐条对照扫描页核对，见进度文件）----
\begin{figure}[H]
\centering
\includegraphics[width=0.60\textwidth]{fig3_2_mu.jpg}
\caption*{\textbf{图 3.2}\ $\Gamma_{t}$（三斜 $P$）的布里渊区。$\Gamma=(000)$；$B=(\frac{1}{2}00)$；$F=(0\frac{1}{2}0)$；$G=(00\frac{1}{2})$。}
\end{figure}
\begin{figure}[H]
\centering
\includegraphics[width=0.64\textwidth]{fig3_3_mu.jpg}
\caption*{\textbf{图 3.3}\ $\Gamma_{m}$（单斜 $P$）的布里渊区。$\Gamma=(000)$；$B=(\bar{\frac{1}{2}}00)$；$Y=(0\frac{1}{2}0)$；$Z=(00\frac{1}{2})$；$C=(0\frac{1}{2}\frac{1}{2})$；$D=(\bar{\frac{1}{2}}0\frac{1}{2})$；$A=(\bar{\frac{1}{2}}\bar{\frac{1}{2}}0)$、$(\bar{\frac{1}{2}}\frac{1}{2}0)$ 或 $(\frac{1}{2}\bar{\frac{1}{2}}0)$；$E=(\bar{\frac{1}{2}}\bar{\frac{1}{2}}\frac{1}{2})$、$(\bar{\frac{1}{2}}\frac{1}{2}\frac{1}{2})$ 或 $(\frac{1}{2}\bar{\frac{1}{2}}\frac{1}{2})$。}
\end{figure}
\begin{figure}[H]
\centering
\includegraphics[width=0.64\textwidth]{fig3_4_mu.jpg}
\caption*{\textbf{图 3.4}\ $\Gamma_{m}^{b}$（单斜 $B$）的布里渊区。$\Gamma=(000)$；$A=(\frac{1}{2}00)$；$Z=(0\frac{1}{2}\frac{1}{2})$；$M=(\frac{1}{2}\frac{1}{2}\frac{1}{2})$；$L=(\frac{1}{2}0\frac{1}{2})$；$V=(00\frac{1}{2})$。}
\end{figure}
\begin{figure}[H]
\centering
\includegraphics[width=0.64\textwidth]{fig3_5_mu.jpg}
\caption*{\textbf{图 3.5}\ $\Gamma_{o}$（正交 $P$）的布里渊区。$\Gamma=(000)$；$Y=(\frac{1}{2}00)$；$X=(0\frac{1}{2}0)$；$Z=(00\frac{1}{2})$；$U=(0\frac{1}{2}\frac{1}{2})$；$T=(\frac{1}{2}0\frac{1}{2})$；$S=(\frac{1}{2}\frac{1}{2}0)$；$R=(\frac{1}{2}\frac{1}{2}\frac{1}{2})$。}
\end{figure}
\begin{figure}[H]
\centering
\includegraphics[width=0.58\textwidth]{fig3_6a_mu.jpg}\\[1ex]
\includegraphics[width=0.58\textwidth]{fig3_6b_mu.jpg}
\caption*{\textbf{图 3.6}\ $\Gamma_{o}^{b}$（正交 $C$）的布里渊区：(a) $a>b$，$\Gamma=(000)$；$Y=(\frac{1}{2}\frac{1}{2}0)$；$Z=(00\frac{1}{2})$；$T=(\frac{1}{2}\frac{1}{2}\frac{1}{2})$；$S=(0\frac{1}{2}0)$；$R=(0\frac{1}{2}\frac{1}{2})$；(b) $b>a$，$\Gamma=(000)$；$Y=(\frac{1}{2}\frac{1}{2}0)$；$Z=(00\frac{1}{2})$；$T=(\frac{1}{2}\frac{1}{2}\frac{1}{2})$；$S=(0\frac{1}{2}0)$；$R=(0\frac{1}{2}\frac{1}{2})$。}
\end{figure}

在继续讲下面的图之前，先说明一下我们的做法。对三斜布拉维格子（图 3.2）和两种单斜布拉维格子（图 3.3 与图 3.4），我们按照定义 3.2.2 用倒空间的初基晶胞来定义布里渊区（见图注 (ii)）。其余格子的布里渊区用倒空间的 Wigner--Seitz 晶胞定义。

\begin{figure}[H]
\centering
\includegraphics[width=0.62\textwidth]{fig3_7ab_mu.jpg}\\[1ex]
\includegraphics[width=0.62\textwidth]{fig3_7c_mu.jpg}
\caption*{\textbf{图 3.7}\ $\Gamma_{o}^{v}$（正交 $I$）的布里渊区：(a) $a>b>c$ 或 $a>c>b$，$\Gamma=(000)$；$X=(\frac{1}{2}\bar{\frac{1}{2}}\frac{1}{2})$；$R=(\bar{\frac{1}{2}}00)$；$S=(\frac{1}{2}0\bar{\frac{1}{2}})$；$T=(\frac{1}{2}\bar{\frac{1}{2}}0)$；$W=(\frac{3}{4}\bar{\frac{1}{2}}\frac{3}{4})$；(b) $b>a>c$ 或 $b>c>a$，$\Gamma=(000)$；$X=(\frac{1}{2}\frac{1}{2}\bar{\frac{1}{2}})$；$R=(\bar{\frac{1}{2}}00)$；$S=(\frac{1}{2}0\bar{\frac{1}{2}})$；$T=(\frac{1}{2}\bar{\frac{1}{2}}0)$；$W=(\frac{3}{4}\bar{\frac{1}{2}}\frac{3}{4})$；(c) $c>b>a$ 或 $c>a>b$，$\Gamma=(000)$；$X=(\frac{1}{2}\bar{\frac{1}{2}}\frac{1}{2})$；$R=(\bar{\frac{1}{2}}00)$；$S=(\frac{1}{2}0\bar{\frac{1}{2}})$；$T=(\frac{1}{2}\bar{\frac{1}{2}}0)$；$W=(\frac{3}{4}\bar{\frac{1}{2}}\frac{3}{4})$。}
\end{figure}
\begin{figure}[H]
\centering
\includegraphics[width=0.60\textwidth]{fig3_8a_mu.jpg}\\[1ex]
\includegraphics[width=0.60\textwidth]{fig3_8b_mu.jpg}\\[1ex]
\includegraphics[width=0.60\textwidth]{fig3_8c_mu.jpg}\\[1ex]
\includegraphics[width=0.60\textwidth]{fig3_8d_mu.jpg}
\caption*{\textbf{图 3.8}\ $\Gamma_{o}^{f}$（正交 $F$）的布里渊区：(a) $1/a^{2}<1/b^{2}+1/c^{2}$、$1/b^{2}<1/c^{2}+1/a^{2}$ 且 $1/c^{2}<1/a^{2}+1/b^{2}$，$\Gamma=(000)$；$Y=(0\bar{\frac{1}{2}}\frac{1}{2})$；$X=(\frac{1}{2}0\frac{1}{2})$；$Z=(\frac{1}{2}\bar{\frac{1}{2}}0)$；$L=(\frac{1}{4}00)$；(b) $1/c^{2}>1/a^{2}+1/b^{2}$，$\Gamma=(000)$；$Y=(0\bar{\frac{1}{2}}\frac{1}{2})$；$X=(\frac{1}{2}0\frac{1}{2})$；$Z=(\frac{1}{2}\bar{\frac{1}{2}}0)$；$L=(\frac{1}{4}00)$；(c) $1/b^{2}>1/a^{2}+1/c^{2}$，$\Gamma=(000)$；$Y=(1\bar{\frac{1}{2}}\frac{1}{2})$；$X=(\frac{1}{2}0\frac{1}{2})$；$Z=(\frac{1}{2}\bar{\frac{1}{2}}0)$；$L=(\frac{1}{4}00)$；(d) $1/a^{2}>1/b^{2}+1/c^{2}$，$\Gamma=(000)$；$Y=(0\bar{\frac{1}{2}}\frac{1}{2})$；$X=(\frac{1}{2}0\frac{1}{2})$；$Z=(\frac{1}{2}\frac{1}{2}0)$；$L=(\frac{1}{4}00)$。}
\end{figure}
\begin{figure}[H]
\centering
\includegraphics[width=0.60\textwidth]{fig3_9_mu.jpg}
\caption*{\textbf{图 3.9}\ $\Gamma_{q}$（四方 $P$）的布里渊区。$\Gamma=(000)$；$M=(\frac{1}{2}\frac{1}{2}0)$；$Z=(00\frac{1}{2})$；$A=(\frac{1}{2}\frac{1}{2}\frac{1}{2})$；$R=(0\frac{1}{2}\frac{1}{2})$；$X=(0\frac{1}{2}0)$。}
\end{figure}
\begin{figure}[H]
\centering
\includegraphics[width=0.60\textwidth]{fig3_10a_mu.jpg}\\[1ex]
\includegraphics[width=0.60\textwidth]{fig3_10b_mu.jpg}
\caption*{\textbf{图 3.10}\ $\Gamma_{q}^{v}$（四方 $I$）的布里渊区：(a) $a>c$；(b) $c>a$。两种情形都有 $\Gamma=(000)$；$N=(0\frac{1}{2}0)$；$X=(00\frac{1}{2})$；$Z=(\bar{\frac{1}{2}}\bar{\frac{1}{2}}\bar{\frac{1}{2}})$；$P=(\frac{1}{4}\frac{1}{4}\frac{1}{4})$。}
\end{figure}
\begin{figure}[H]
\centering
\includegraphics[width=0.60\textwidth]{fig3_11a_mu.jpg}\\[1ex]
\includegraphics[width=0.60\textwidth]{fig3_11b_mu.jpg}
\caption*{\textbf{图 3.11}\ $\Gamma_{rh}$（三方 $R$）的布里渊区：(a) $a>\sqrt{2}\,c$，$\Gamma=(000)$；$Z=(\frac{1}{2}\frac{1}{2}\bar{\frac{1}{2}})$；$L=(0\frac{1}{2}0)$；$F=(0\frac{1}{2}\bar{\frac{1}{2}})$；(b) $\sqrt{2}\,c>a$，$\Gamma=(000)$；$Z=(\frac{1}{2}\frac{1}{2}\frac{1}{2})$；$L=(0\frac{1}{2}0)$；$F=(\frac{1}{2}\frac{1}{2}0)$。在 (a) 中 $\Sigma$ 和 $Y$ 是两条完整的对称线；在 (b) 中 $Q$ 和 $B$ 经过一个三次转动分别与 $Q'$ 和 $B'$ 相关联，二者分别是线 $\Sigma$ 和 $Y$ 的延续（$F'=(0\bar{\frac{1}{2}}\bar{\frac{1}{2}})$，由 $F$ 经三次转动得到；另见表 3.6）。}
\end{figure}
\begin{figure}[H]
\centering
\includegraphics[width=0.60\textwidth]{fig3_12_mu.jpg}
\caption*{\textbf{图 3.12}\ $\Gamma_{h}$（六角 $P$）的布里渊区。$\Gamma=(000)$；$M=(0\frac{1}{2}0)$；$A=(00\frac{1}{2})$；$L=(0\frac{1}{2}\frac{1}{2})$；$K=(\bar{\frac{1}{3}}\frac{2}{3}0)$；$H=(\bar{\frac{1}{3}}\frac{2}{3}\frac{1}{2})$。}
\end{figure}
\begin{figure}[H]
\centering
\includegraphics[width=0.58\textwidth]{fig3_13_mu.jpg}
\caption*{\textbf{图 3.13}\ $\Gamma_{c}$（立方 $P$）的布里渊区。$\Gamma=(000)$；$X=(0\frac{1}{2}0)$；$M=(\frac{1}{2}\frac{1}{2}0)$；$R=(\frac{1}{2}\frac{1}{2}\frac{1}{2})$。}
\end{figure}
\begin{figure}[H]
\centering
\includegraphics[width=0.58\textwidth]{fig3_14_mu.jpg}
\caption*{\textbf{图 3.14}\ $\Gamma_{c}^{f}$（立方 $F$）的布里渊区。$\Gamma=(000)$；$X=(\frac{1}{2}0\frac{1}{2})$；$L=(\frac{1}{2}\frac{1}{2}\frac{1}{2})$；$W=(\frac{1}{2}\frac{1}{4}\frac{3}{4})$。$K$ 和 $U$ 虽常被称为对称点，但按定义 3.3.1 它们并不算对称点，因为它们并不比相邻的线 $\Sigma$ 和 $S$ 具有更多对称性。}
\end{figure}
\begin{figure}[H]
\centering
\includegraphics[width=0.56\textwidth]{fig3_15_mu.jpg}
\caption*{\textbf{图 3.15}\ $\Gamma_{c}^{v}$（立方 $I$）的布里渊区。$\Gamma=(000)$；$H=(\frac{1}{2}\frac{1}{2}\frac{1}{2})$；$P=(\frac{1}{4}\frac{1}{4}\frac{1}{4})$；$N=(00\frac{1}{2})$。}
\end{figure}

对某些布拉维格子，布里渊区的形状是唯一的，例如三种立方布拉维格子（其布里渊区绘于图 3.13--3.15）。但对另外一些布拉维格子，可能出现的形状不止一种，取决于布拉维格子的轴比和轴间角。对三斜布拉维格子（图 3.2）和两种单斜布拉维格子（图 3.3、3.4），我们遵循定义 3.2.2，用倒空间的初基晶胞来定义布里渊区。对另一些布拉维格子，可能出现表面上不同的布里渊区（例如图 3.8 的 $\Gamma_{o}^{f}$、图 3.10 的 $\Gamma_{q}^{v}$、图 3.11 的 $\Gamma_{rh}$）。然而，当可能有两个表面上不同的布里渊区时，每种情形中对称点的数目是相同的。例如在图 3.11(a)（$\Gamma_{rh}$，$a>\sqrt{2}\,c$）中，对称点是 $\Gamma$（$\mathbf{k}=0$）、$L$（$\mathbf{k}=\frac{1}{2}\mathbf{g}_{2}$）、$F$（$\mathbf{k}=\frac{1}{2}\mathbf{g}_{2}-\frac{1}{2}\mathbf{g}_{3}$）和 $Z$（$\mathbf{k}=\frac{1}{2}\mathbf{g}_{1}+\frac{1}{2}\mathbf{g}_{2}-\frac{1}{2}\mathbf{g}_{3}$）；而在图 3.11(b) 的替代布里渊区（$a<\sqrt{2}\,c$）中，对称点是 $\Gamma$（$\mathbf{k}=0$）、$L$（$\mathbf{k}=\frac{1}{2}\mathbf{g}_{2}$）、$F$（$\mathbf{k}=\frac{1}{2}\mathbf{g}_{1}+\frac{1}{2}\mathbf{g}_{2}$）和 $Z$（$\mathbf{k}=\frac{1}{2}\mathbf{g}_{1}+\frac{1}{2}\mathbf{g}_{2}+\frac{1}{2}\mathbf{g}_{3}$）。$\Gamma$ 和 $L$ 在两个布里渊区中以相同的 $\mathbf{k}$ 出现，而图 3.11(a) 中的 $F$ 和 $Z$ 在图 3.11(b) 中对应的 $\mathbf{k}$ 矢量不同但等价。对另外一些布拉维格子，布里渊区的外形相同，只是各面相对 $k_{x}$、$k_{y}$、$k_{z}$ 轴的取向不同，见图 3.16。

\begin{figure}[H]
\centering
\includegraphics[width=0.56\textwidth]{fig3_16_mu.jpg}
\caption*{\textbf{图 3.16}\ $\Gamma_{o}^{b}$ 的布里渊区。}
\end{figure}

考虑 $\Gamma_{o}^{b}$（图 3.6）。当 $a>b$ 时（图 3.6(a)），设 $T$ 是线 $ZT$ 上的一个点，其 $\mathbf{k}$ 矢量为 $\mathbf{k}_{T}$，则
\begin{equation*}
\mathbf{k}_{E} = \mathbf{k}_{T} + (\alpha\mathbf{g}_{1} + \alpha\mathbf{g}_{2}).
\tag{3.3.2}
\end{equation*}
图 3.6(a) 中相应的对称线具有由式 (3.3.2) 给出的 $\mathbf{k}$ 矢量 $\mathbf{k}_{E}$，即线 $ZT$ 的延续；见图 3.16。这条线 $E$ 现在已跑到第一布里渊区之外，但它等价于第一区内的一条线，即 $E'$，其中
\begin{equation*}
\mathbf{k}_{E'} = \mathbf{k}_{E} - \mathbf{g}_{1} - \mathbf{g}_{2}.
\tag{3.3.3}
\end{equation*}
线 $E'$ 与图 3.6(a) 中的线 $A$ 有简单的对称关系，因此不必把它同 $A$ 分开考虑。反过来看，这意味着图 3.6(a) 中的线 $A$ 在图 3.6(b) 中被分成两段：对某些 $\alpha$ 值的一段是图 3.6(a) 的 $A$，对其余 $\alpha$ 值的一段是图 3.6(b) 的 $E$。类似地，图 3.6(a) 中的线 $\Sigma$ 在图 3.6(b) 中被分成 $\Lambda$ 和 $C$。反过来，图 3.6(b) 中的线 $A$ 在图 3.6(a) 中被分成 $\Delta$ 和 $F$，而图 3.6(b) 中的线 $B$ 被分成 $\Delta$ 和 $F'$。

由于我们指出的这种等价性，在本书后面的表中，对 14 种布拉维格子中的每一种，只需对一个布里渊区列出对称点即可；通常也只需对每一种格子列出一张对称线表。不过，对区面上的对称线存在少数例外。例如，虽然图 3.6(b) 中的 $C$ 与图 3.6(a) 中的 $E$ 相关联，但二者之间存在一个微妙的差别，这将在后文变得明显。

14 种布拉维格子各自基本域中每个对称点和对称线的群 $\mathbf{P}(\mathbf{k})$（$\mathbf{k}$ 的对称群）列于表 3.6。可以注意到，图 3.2--3.15 中基本域的许多棱并不是对称线，而只是（按定义 3.3.2 的意义）对称面。把我们限制在单一基本域的理由后文自明；例如见 5.5 节。

\begin{figure}[H]
\centering
\includegraphics[width=0.92\textwidth]{c3_t36_p1.png}
\caption*{\textbf{表 3.6}\ 对称群 $\mathbf{P}(\mathbf{k})$（原书第 113 页起扫描；第 1--5 页）。}
\end{figure}
\begin{figure}[H]
\centering
\includegraphics[width=0.92\textwidth]{c3_t36_p2.png}
\caption*{\textbf{表 3.6}（第 2 页）。}
\end{figure}
\begin{figure}[H]
\centering
\includegraphics[width=0.92\textwidth]{c3_t36_p3.png}
\caption*{\textbf{表 3.6}（第 3 页）。}
\end{figure}
\begin{figure}[H]
\centering
\includegraphics[width=0.92\textwidth]{c3_t36_p4.png}
\caption*{\textbf{表 3.6}（第 4 页）。}
\end{figure}
\begin{figure}[H]
\centering
\includegraphics[width=0.92\textwidth]{c3_t36_p5.png}
\caption*{\textbf{表 3.6}（第 5 页；本页下部有关于 $\Gamma_{rh}$ 的 $F$ 点的脚注）。}
\end{figure}
\begin{figure}[H]
\centering
\includegraphics[width=0.92\textwidth]{c3_t36_p6.png}
\caption*{\textbf{表 3.6}（第 6 页，立方 $P$、$F$、$I$；页末起为表的注）。}
\end{figure}
\begin{figure}[H]
\centering
\includegraphics[width=0.92\textwidth]{c3_t36_p7.png}
\caption*{\textbf{表 3.6}（注 (ii)--(iv)，原书扫描占位）。}
\end{figure}

表 3.6 的注（译文）：

\noindent（i）第 1 列列出正空间的布拉维格子，并指出图 3.2--3.15 中展示各对称点和对称线位置的相应部分。

\noindent（ii）第 2 列列出对称点和对称线。它们既可从图 3.2--3.15 辨认，也可从第 3 列辨认；第 3 列列出各点和线相对于 $\mathbf{g}_{1}$、$\mathbf{g}_{2}$、$\mathbf{g}_{3}$ 的坐标。例如 $\Gamma_{m}$ 中的 $W$ 是线 $YC$ 上的一个点，其 $\mathbf{k}$ 矢量为 $\frac{1}{2}\mathbf{g}_{2}+\alpha\mathbf{g}_{3}$（$0<\alpha<\frac{1}{2}$）。对称点和对称线所用的标记，除少数例外，都是 Miller and Love（1967）的记号。

\noindent（iii）第 4 列列出各点和线的对称群 $\mathbf{P}(\mathbf{k})$ 的名称。例如 $\Gamma_{m}$ 中的 $W$ 具有对称群 $2\ (C_{2})$。对称群的定义见 3.3 节正文。矢量 $\mathbf{g}_{1}$、$\mathbf{g}_{2}$、$\mathbf{g}_{3}$ 相对于笛卡尔轴 $\Gamma k_{x}k_{y}k_{z}$ 的取向见表 3.3。

\noindent（iv）第 5 列列出对称群 $\mathbf{P}(\mathbf{k})$ 的元素。这些元素可借助图 1.1--1.3 来辨认，此时在所有情形下都取 $Oxyz$ 轴与图 3.2--3.15 的 $\Gamma k_{x}k_{y}k_{z}$ 轴重合。

\noindent（关于表中 $\Gamma_{rh}$ 的 $F$ 点的脚注：为了把 $F$ 保持在所选的基本域内，不得不让 (b) 中 $F$ 的 $\mathbf{P}(\mathbf{k})$ 元素与 (a) 中 $F$ 的不同，尽管这两个点显然是互相关联的。若允许 $F$ 位于所选基本域之外（即位于图 3.11(b) 的 $F'$ 处），本可使二者一致，但我们不拟这样做。）
